% Section 10.2, from lectures/L10/appendix/b_device_tree.md.

\section{The Device Tree}
\label{c10:app:b}

Where the description of the hardware comes from, how to read one, and how a driver gets its
addresses out of it without containing any.

\S\ref{c6:app:b1} did the translation by hand once. This is the rest of the format, and the
machinery that does it for you.

\subsection{Why it exists}
\label{c10:app:b1}

An x86 PC can discover itself. PCI is enumerable, ACPI describes the rest, and one kernel image
boots on any machine.

An SoC can do none of that. There is no bus to walk that reveals a UART at \code{0x09000000} raising
interrupt 33. Somebody has to say so, and before 2011 that somebody was C code in \file{arch/arm/},
one file per board, thousands of them, all of which had to be compiled into the kernel and
maintained by the kernel community. The device tree moved that description out of the kernel and
into a data file.

The consequence is the one that matters commercially: \textbf{one kernel image, many boards.} The
bootloader passes a different DTB and the same binary boots different hardware.

\subsection{The syntax}
\label{c10:app:b2}

\begin{console}
/ {
        #address-cells = <2>;
        #size-cells = <2>;

        platform-bus@c000000 {
                compatible = "qemu,platform", "simple-bus";
                #address-cells = <1>;
                #size-cells = <1>;
                ranges = <0x00 0x00 0xc000000 0x2000000>;
                interrupt-parent = <0x8003>;

                qa-dev@0 {
                        compatible = "qacademy,qa-dev-1.0";
                        reg = <0x00 0x1000>;
                        interrupts = <0x00 0x70 0x04>;
                };
        };
};
\end{console}

\begin{booktable}{@{}lL@{}}
  \thead{Element} & \thead{Is} \\
  \midrule
  \code{name@address} & A node. The part after \code{@} matches the first address in \code{reg} \\
  \code{property = value} & A property. Values are byte strings, printed as cells or text \\
  \code{<...>} & Cells: 32-bit big-endian words \\
  \code{"..."} & A string, or a list of them separated by nulls \\
  \code{<&label>} & A phandle: a pointer to another node \\
\end{booktable}

\noindent Node names are conventional and carry no meaning to the kernel; \textbf{\code{compatible}
is what matters.}

\subsection{\code{compatible}}
\label{c10:app:b3}

The match key, and a list rather than a string:

\begin{console}
compatible = "qemu,platform", "simple-bus";
\end{console}

\noindent \textbf{Most specific first.} A driver matches if its \code{of_match_table} contains any
of the strings, and where its table contains several, the entry for the earliest string wins. So a
node can say ``I am exactly this chip, and failing that I behave like this generic thing'', and a
kernel with only the generic driver still works. The order does not rank one driver above another:
if a driver for each string is loaded, whichever the bus tries first binds.

The convention is \code{vendor,model}. The vendor prefix is registered in
\file{Documentation/devicetree/bindings/vendor-prefixes.yaml}, and making one up is how a binding
gets rejected upstream.

\textbf{A driver matches on the string, not on the node name.} Change \code{qa-dev@0} to
\code{widget@0} and everything still works; change one character of \code{qacademy,qa-dev-1.0} and
\code{probe} is never called. That is the most common reason a driver silently does not load, and
the diagnosis is to compare the two strings by eye:

\begin{console}
# tr -d '\0' < /sys/firmware/devicetree/base/platform-bus@c000000/qa-dev@0/compatible
qacademy,qa-dev-1.0
\end{console}

\subsection{\code{reg}, cells, and \code{ranges}}
\label{c10:app:b4}

\code{reg} is a list of address-and-size pairs, but \textbf{how many cells each takes is declared by
the parent}, not by the node itself:

\begin{console}
parent:  #address-cells = <1>;  #size-cells = <1>;
child:   reg = <0x00 0x1000>;         -> one region: address 0x0, size 0x1000
\end{console}

\noindent With \code{\#address-cells = <2>} the same property would be one address of two cells and
no size, or an error. \textbf{You cannot read a \code{reg} without looking at its parent}, and that
is the single most common mistake in reading a device tree by hand.

\code{ranges} translates a child address into the parent's address space, as a list of
\code{<child-address parent-address size>}:

\begin{console}
ranges = <0x00 0x00 0xc000000 0x2000000>;
\end{console}

\noindent Here the child address takes 1 cell, the parent address 2 (the root declares
\code{\#address-cells = <2>}), and the size 1. So: child \code{0x0} maps to parent
\code{0x0_0c000000}, for \code{0x2000000} bytes. The device at child offset \code{0x0} is therefore
at physical \textbf{\code{0xC000000}}.

An empty \code{ranges;} means the child addresses \emph{are} the parent's, with no translation. No
\code{ranges} property at all means the node is not memory-mapped through its parent, and children
cannot be translated.

\textbf{The cells are big-endian, always}, regardless of the machine. Reading them with \code{od
-tx4} on a little-endian host byte-swaps every one, so \code{0x0C000000} reads as \code{0x0000000C},
which is plausible enough to act on. Use \code{-tx1} and reassemble, or use a tool that knows.

\subsection{Interrupts}
\label{c10:app:b5}

\begin{console}
interrupts = <0x00 0x70 0x04>;
interrupt-parent = <0x8003>;
\end{console}

\noindent \code{interrupts} is interpreted \textbf{by the interrupt controller}, not by the kernel
generally, and how many cells it takes and what they mean is that controller's business. For a GIC
it is three:

\begin{booktable}{@{}r>{\hsize=1.29\hsize}L>{\hsize=0.71\hsize}L@{}}
  \thead{Cell} & \thead{Meaning} & \thead{Here} \\
  \midrule
  0 & Type: 0 = SPI, 1 = PPI & 0, a shared interrupt \\
  1 & Number within that type & 0x70 = 112 \\
  2 & Flags: 1 = rising edge, 4 = level high & 4, level triggered \\
\end{booktable}

\noindent The GIC's own numbering adds the 32-entry SPI base, so this is hardware interrupt
\textbf{144}, and Linux allocates its own number on top of that. All three appear in one line of
\file{/proc/interrupts}, which \S\ref{c8:app:a1} takes apart.

\code{interrupt-parent} is a phandle naming which controller, inherited by children if not
overridden.

\subsection{Getting it into your driver}
\label{c10:app:b6}

The point of the whole format:

\begin{ccode}
priv->regs = devm_platform_ioremap_resource(pdev, 0);
priv->irq  = platform_get_irq(pdev, 0);
\end{ccode}

\noindent \code{devm_platform_ioremap_resource} does, in one call, what \S\ref{c6:app:b7} spent a
page on: takes \code{reg} entry 0, which the kernel already translated through every \code{ranges}
between the node and the root when it created the device, claims the region, maps it, and registers
a destructor for all of it.

\code{platform_get_irq} parses \code{interrupts}, finds the controller through
\code{interrupt-parent}, creates the mapping and returns the Linux number. \textbf{It returns a
negative error, not zero, on failure}, so it is checked with \code{< 0}. Treating 0 as the failure
value is a bug that never shows on a board where the call succeeds, and on the first board where it
fails, passes the negative error on to \code{devm_request_irq} as if it were a number, deferral
included.

Other properties are read with the \code{of_property_*} family:

\begin{ccode}
u32 depth;
if (of_property_read_u32(pdev->dev.of_node, "fifo-depth", &depth))
        depth = 16;                     /* absent: use the default */
\end{ccode}

\noindent Note the idiom: a missing optional property is not an error, it is a default. A missing
\emph{required} property is an error and should say which property, by name, in \code{dev_err}.

\subsection{Bindings}
\label{c10:app:b7}

A \textbf{binding} is the documentation of what properties a \code{compatible} string implies. They
live in \file{Documentation/devicetree/bindings/}, and since about 2018 they are YAML schemas rather
than prose, which means they can be checked mechanically:

\begin{shell}
make dt_binding_check          # are the schemas themselves valid
make dtbs_check                # do the shipped device trees match them
\end{shell}

\noindent This matters more than it sounds. \textbf{A device tree is an ABI.} A DTB written for one
kernel is expected to work on later ones, because the bootloader that supplies it may not be updated
with the kernel. So a binding, once accepted, is a promise, and getting one reviewed upstream is
deliberately harder than getting a driver reviewed.

Two consequences for a driver author:
\begin{itemize}
  \item \textbf{Do not invent properties for things that are not hardware description.} A debug flag
    or a buffer size that is a policy choice is a module parameter, not a device tree property.
  \item \textbf{Read the binding before the driver.} For an unfamiliar device the binding tells you
    what the hardware is; the driver tells you only what one author did about it.
\end{itemize}

\subsection{Where the DTB comes from, and overlays}
\label{c10:app:b8}

At boot, the bootloader loads the DTB into memory and passes its address to the kernel in a
register. \S\ref{c1:app:a3} drew that handover; this is the file it was handing over.

On this course's target there is no bootloader: \textbf{QEMU generates the DTB itself}, including
the \code{qa-dev} node, and hands it to the kernel exactly as U-Boot would. The node's address and
interrupt are whatever QEMU assigned when it placed the device, which is why nobody typed them
anywhere. You can dump what it produced:

\begin{shell}
qemu-system-aarch64 -machine virt,dumpdtb=virt.dtb ... -device qa-dev
dtc -I dtb -O dts virt.dtb
\end{shell}

\noindent \textbf{Overlays} are fragments applied to a base DTB at run time, which is how an
expansion board or a cape describes itself without a new base tree. They are genuinely useful for
that and are commonly reached for as a way to hand-patch a tree during development, which works and
is worth knowing is a development technique rather than a shipping one.

\subsection{Reading one on a running system}
\label{c10:app:b9}

The whole tree is under \file{/sys/firmware/devicetree/base}, one directory per node and one file
per property:

\begin{console}
/sys/firmware/devicetree/base/platform-bus@c000000/qa-dev@0/
    compatible      "qacademy,qa-dev-1.0"
    reg             00 00 00 00 00 00 10 00
    interrupts      00 00 00 00 00 00 00 70 00 00 00 04
\end{console}

\noindent Strings come out with trailing nulls, so \code{tr -d '\\0'} is usually wanted. Numeric
properties are raw big-endian bytes and need reassembling by hand, per \S\ref{c10:app:b4}.

The other view is \file{/sys/bus/platform/devices/}, where each node with a \code{compatible} has
become a device named from its translated address. Those two views are the same information from the
two sides of \S\ref{c10:app:b1}: what the firmware said, and what the kernel made of it.
